\documentclass[11pt]{article}

\usepackage[margin=1in]{geometry}
\usepackage[T1]{fontenc}
\usepackage{lmodern}
\usepackage{microtype}
\usepackage{amsmath}
\usepackage{booktabs}
\usepackage{graphicx}
\usepackage[hidelinks]{hyperref}

\title{Short-Horizon Forecasting of the SPY Implied-Volatility Surface\\
  \large A pre-registered study on Databento OPRA data}
\author{Viranca}
\date{\today}

\begin{document}
\maketitle

\begin{abstract}
% One paragraph: the question, the data, the headline result (including what
% did NOT work), and why the methodology (pre-registration, execution-aware
% metrics) makes the result trustworthy.
TODO.
\end{abstract}

\section{Question and setup}
% What is being predicted (surface state deltas at 15m/60m horizons, short-DTE
% wings), why short-horizon surface dynamics matter, and what "good" means
% here (beating persistence after spreads, not raw RMSE).

\section{Data and surface construction}
% Databento OPRA CMBP-1 for SPY, one month; surface-state builder;
% arbitrage/quality gates; reference-IV cross-checks.

\section{Methodology and pre-registration}
% Baselines (persistence), models (ridge, LightGBM), walk-forward splits,
% expectations registered before each run, and the evaluation metrics.

\section{Results}
% Headline table: model vs. persistence by horizon. Feature-set ablations
% (Greeks, microstructure). Report negative results as registered.

\section{What changed between rounds --- and why}
% The replication story: which June findings held on three months, which
% flipped sign, and the diagnosis.

\section{Limitations and next steps}
% Single underlier, one regime, execution model is a thin spread-aware
% wrapper; what a production version would add.

\end{document}
